
%%%%%%%%%%%%%%%%%%%%%%%%%%%%%%%%%%%%%%%%%%%%%%%%%%%%%%%%%%%%%%%%%%%%%%%%%%%%%%

%\section{Preliminaries}
%\label{sec: preliminaries}
%
%%%%%%%%%%%%%%%%%%%%%%%%%%%%%%%%%%%%%%%%%%%%%%%%%%%%%%%%%%%%%%%%%%%%%%%%%%%%%%

%\section{Preliminaries}
%\label{sec: preliminaries}
%
%%%%%%%%%%%%%%%%%%%%%%%%%%%%%%%%%%%%%%%%%%%%%%%%%%%%%%%%%%%%%%%%%%%%%%%%%%%%%%

%\section{Preliminaries}
%\label{sec: preliminaries}
%
%%%%%%%%%%%%%%%%%%%%%%%%%%%%%%%%%%%%%%%%%%%%%%%%%%%%%%%%%%%%%%%%%%%%%%%%%%%%%%

%\section{Preliminaries}
%\label{sec: preliminaries}
%\input{2_preliminaries.tex}

Let $\binrationals$ be the set of  binary numbers.
For $c\in\binrationals$, we denote by $\precision(c)$ the precision of $c$, that is the
number of fractional bits in the binary representation of $c$, assuming there are no trailing $0$'s.
In other words, $\precision(c)$ is the smallest $d\in\nonnegintegers$ such that $c = a/2^d$ for some $a\in\integers$. 
If $c = a/2^d$ represents actual fluid concentration, then we have $0\le a \le 2^d$. However,  
it is convenient to relax this restriction and allow ``concentration values'' that are
arbitrary binary numbers, even negative. In fact, as we show shortly, 
it will be convenient to work with integral values.

By a \emph{configuration} we mean a multiset of $n$ binary numbers, called droplets or concentrations.
In the literature, multisets are often represented by their characteristic functions (that specify
the multiplicity of each element of the ground set). In this paper we will generally use set-theoretic
terminology, with its natural interpretation. For example, for a configuration $C$ and a concentration $a$, 
$a\in C$ means that the multiplicity of $a$ in $C$ is strictly positive, while $a\notin C$ means that it's zero.
Or, $C - \braced{a} = C'$ means that the multiplicity of $a$ in $C'$ is one less than in $C$, while
other multiplicities are the same.
The number of droplets in $C$ is denoted $\nofdroplets{C}=n$, 
while the number of different concentrations is denoted $\nofconcentrations{C}=m$.
We will typically denote a configuration by
$C = \braced{f_1:c_1,f_2: c_2,...,f_m:c_m} \ingroundset \binrationals$, where each $c_i$ represents
a (different) concentration value and $f_i$ denotes the
multiplicity of $c_i$ in $C$, so that $\sum_{i=1}^m f_i = n$. 
Occasionally, if it does not lead to confusion, we may say
``droplet $c_i$'' or ``concentration $c_i$'', referring to some droplet with concentration $c_i$.
If $f_i=1$, we shorten ``$f_i:c_i$'' to just ``$c_i$''.
If $f_i = 1$ we say that droplet $c_i$ is a \emph{singleton}, 
if $f_i = 2$ we say that droplet $c_i$ is a \emph{doubleton}
and if $f_i\geq 2$ we say that droplet $c_i$ is a \emph{non-singleton}.
By $\mysum(C)$ we denote the sum of $C$, that is $\mysum(C) = \sum_{c\in C} c$.
$\average(C) = \mysum(C)/n$ is the average value of the concentrations in $C$ and will
be typically denoted by $\mu$.
(Later, we will typically deal with configurations $C$ such that
$C\cup\braced{\mu}\ingroundset\integers$.)

Mixing graphs were defined in the introduction. As we are not concerned in this paper with the
topological properties of mixing graphs, we will often identify a mixing graph $G$ with a
corresponding \emph{mixing sequence}, which is a sequence (not necessarily unique) of mixing operations
that convert $C$ into its perfect mixture. In other words, a mixing sequence is a sequence
of mixing operations in a topological ordering of a mixing graph.

Of course in a perfect-mixing graph (or sequence) $G$ for $C$, all concentrations in $G$,
including those in $C\cup\braced{\mu}$,
must have finite precision (that is, belong to $\binrationals$);
in fact, the maximum concentration in $G$ is at least $\max\braced{\precision(C),\precision(\mu)}$.
In addition to the basic question about finding a perfect-mixing graph for $C$, 
we are also interested in bounding the precision required to do so.

For $x\in\binrationals$, define multisets $C+x = \braced{c+x \mid c\in C}$,
$C-x = C+(-x)$, and $C\cdot x = \braced{c\cdot x \mid c\in C}$.
The next observation says that offsetting all values in $C$
does not affect perfect mixability, as long as the offset value's precision
does not exceed that of $C$ or $\mu$.

\begin{observation}%--------
\label{obs: offsetting does not affect reachability}
Let $\mu = \average(C)$ and $x\in\binrationals$. 
Also, let $d\in\posintegers$ be such that $d\ge \max\braced{\precision(C),\precision(\mu),\precision(x)}$.
Then $C$ is perfectly mixable with precision $d$ if and only if
$C' = C+x$ is perfectly mixable with precision $d$.
\end{observation}%----------

\begin{proof}
$(\Rightarrow)$
Suppose that $G$ is a perfect-mixing sequence for $C$ with precision $d$.
Run the same sequence $G$ on input $C'$. If some mixing step in $G$ produces a value $z$ when
the input is $C$ then on $C'$ its value is $z+x$,
and $\precision(z+x)\le \max\braced{\precision(z),\precision(x)} \le d$.
Thus the maximum precision in $G$ for input $C'$ is at most $d$. 

$(\Leftarrow)$
The proof for this implication follows from noting that
$\mu' = \average(C') = \mu+x$,
$\max\braced{\precision(C'),\precision(\mu')}\le d$, and by applying the
above argument to $-x$ instead of $x$.
\end{proof}

\begin{observation}%--------
\label{obs: power of 2 rescaling}
Let $\mu = \average(C)$, $\delta = \max\braced{\precision(C),\precision(\mu)}$, 
$C' = C\cdot 2^\delta$ with $\mu'=\average(C')= 2^\delta\mu$. 
	(Thus $C'\cup\braced{\mu'}\ingroundset\integers$.)
	Then $C$ is perfectly mixable with precision $d\ge\delta$
	if and only if $C'$ is perfectly mixable with precision $d' = d-\delta$. 
\end{observation}%----------


\begin{proof}
$(\Rightarrow)$
Let $G$ be a perfect-mixing sequence for $C$, with precision $d$.
Run the same sequence $G$ on input $C'$.
If some node in $G$ produces a value $z$ on input $C$, then its value on input $C'$ will be
$z 2^\delta$, and $\precision(z 2^\delta) = \max\braced{\precision(z)-\delta,0} \le d-\delta = d'$.

$(\Leftarrow)$
Let $G'$ be a perfect-mixing sequence for $C'$ with precision $d'$.
Run $G'$ on input $C$. If some node in $G'$ produces a value $y$ on input $C'$, then
its value on input $C$ will be $y/2^\delta$, and
$\precision(y/2^\delta) \leq \precision(y)+\delta \le d'+\delta = d$.
\end{proof}


%%%%%%%

\myparagraph{Integral configurations}
Per Observation~\ref{obs: power of 2 rescaling}, we can restrict our attention to configurations with integer
values and average, that is, we will be assuming that $C\cup\braced{\mu}\ingroundset \integers$.

For $x\in\posintegers$, if each $c\in C$ is a multiple of $x$,
let $C/x = \braced{c/x \mid c\in C}$.
For integral configurations, we can extend Observation~\ref{obs: power of 2 rescaling}
to also multiplying $C$ by an odd integer or dividing it by a common odd factor of
all concentrations in $C$. 

\begin{observation}%--------
\label{obs: odd integer rescaling}
Assume that $C\cup\braced{\mu}\ingroundset \integers$ and let $x\in\posintegers$ be odd.
%
\begin{description}
%
\item{\emph{(a)}} Let $C' = C\cdot x$. Then
	$C$ is perfectly mixable with precision $0$ if and only if
	$C'$ is perfectly mixable with precision $0$.	
%
\item{\emph{(b)}} Suppose that $x$ 
is a divisor of all concentrations in $C\cup\braced{\mu}$.
Then $C$ is perfectly mixable with precision $0$ if and only if
$C/x$ is perfectly mixable with precision $0$.
\end{description}	
\end{observation}%----------

\begin{proof}
Part~(b) follows from~(a), so we only prove part~(a).
Any sequence $G$ of mixing operations for $C$ can be applied to $C\cdot x$. 
By simple induction,
if some intermediate value in $G$ was an integer $z$, now its value will be $z x$,
also an integer. This shows the $(\Rightarrow)$ implication.

To justify the $(\Leftarrow)$ implication, 
suppose that $G'$ is a sequence of mixing operations for $C'$ and that all
concentrations in $G'$ are integer. 
Since $x$ is odd and all concentrations in $C'$ are multiples of $x$, 
every concentration in $G'$ will be also a multiple of $x$, including  $\average(C')$.  
Thus, if  we run $G'$ on $C$ instead of $C'$, if some node's concentration
was $c x$, now it will be $c$.
Thus, the $(\Leftarrow)$ implication holds.
\end{proof}

%%%%%%%








Let $\binrationals$ be the set of  binary numbers.
For $c\in\binrationals$, we denote by $\precision(c)$ the precision of $c$, that is the
number of fractional bits in the binary representation of $c$, assuming there are no trailing $0$'s.
In other words, $\precision(c)$ is the smallest $d\in\nonnegintegers$ such that $c = a/2^d$ for some $a\in\integers$. 
If $c = a/2^d$ represents actual fluid concentration, then we have $0\le a \le 2^d$. However,  
it is convenient to relax this restriction and allow ``concentration values'' that are
arbitrary binary numbers, even negative. In fact, as we show shortly, 
it will be convenient to work with integral values.

By a \emph{configuration} we mean a multiset of $n$ binary numbers, called droplets or concentrations.
In the literature, multisets are often represented by their characteristic functions (that specify
the multiplicity of each element of the ground set). In this paper we will generally use set-theoretic
terminology, with its natural interpretation. For example, for a configuration $C$ and a concentration $a$, 
$a\in C$ means that the multiplicity of $a$ in $C$ is strictly positive, while $a\notin C$ means that it's zero.
Or, $C - \braced{a} = C'$ means that the multiplicity of $a$ in $C'$ is one less than in $C$, while
other multiplicities are the same.
The number of droplets in $C$ is denoted $\nofdroplets{C}=n$, 
while the number of different concentrations is denoted $\nofconcentrations{C}=m$.
We will typically denote a configuration by
$C = \braced{f_1:c_1,f_2: c_2,...,f_m:c_m} \ingroundset \binrationals$, where each $c_i$ represents
a (different) concentration value and $f_i$ denotes the
multiplicity of $c_i$ in $C$, so that $\sum_{i=1}^m f_i = n$. 
Occasionally, if it does not lead to confusion, we may say
``droplet $c_i$'' or ``concentration $c_i$'', referring to some droplet with concentration $c_i$.
If $f_i=1$, we shorten ``$f_i:c_i$'' to just ``$c_i$''.
If $f_i = 1$ we say that droplet $c_i$ is a \emph{singleton}, 
if $f_i = 2$ we say that droplet $c_i$ is a \emph{doubleton}
and if $f_i\geq 2$ we say that droplet $c_i$ is a \emph{non-singleton}.
By $\mysum(C)$ we denote the sum of $C$, that is $\mysum(C) = \sum_{c\in C} c$.
$\average(C) = \mysum(C)/n$ is the average value of the concentrations in $C$ and will
be typically denoted by $\mu$.
(Later, we will typically deal with configurations $C$ such that
$C\cup\braced{\mu}\ingroundset\integers$.)

Mixing graphs were defined in the introduction. As we are not concerned in this paper with the
topological properties of mixing graphs, we will often identify a mixing graph $G$ with a
corresponding \emph{mixing sequence}, which is a sequence (not necessarily unique) of mixing operations
that convert $C$ into its perfect mixture. In other words, a mixing sequence is a sequence
of mixing operations in a topological ordering of a mixing graph.

Of course in a perfect-mixing graph (or sequence) $G$ for $C$, all concentrations in $G$,
including those in $C\cup\braced{\mu}$,
must have finite precision (that is, belong to $\binrationals$);
in fact, the maximum concentration in $G$ is at least $\max\braced{\precision(C),\precision(\mu)}$.
In addition to the basic question about finding a perfect-mixing graph for $C$, 
we are also interested in bounding the precision required to do so.

For $x\in\binrationals$, define multisets $C+x = \braced{c+x \mid c\in C}$,
$C-x = C+(-x)$, and $C\cdot x = \braced{c\cdot x \mid c\in C}$.
The next observation says that offsetting all values in $C$
does not affect perfect mixability, as long as the offset value's precision
does not exceed that of $C$ or $\mu$.

\begin{observation}%--------
\label{obs: offsetting does not affect reachability}
Let $\mu = \average(C)$ and $x\in\binrationals$. 
Also, let $d\in\posintegers$ be such that $d\ge \max\braced{\precision(C),\precision(\mu),\precision(x)}$.
Then $C$ is perfectly mixable with precision $d$ if and only if
$C' = C+x$ is perfectly mixable with precision $d$.
\end{observation}%----------

\begin{proof}
$(\Rightarrow)$
Suppose that $G$ is a perfect-mixing sequence for $C$ with precision $d$.
Run the same sequence $G$ on input $C'$. If some mixing step in $G$ produces a value $z$ when
the input is $C$ then on $C'$ its value is $z+x$,
and $\precision(z+x)\le \max\braced{\precision(z),\precision(x)} \le d$.
Thus the maximum precision in $G$ for input $C'$ is at most $d$. 

$(\Leftarrow)$
The proof for this implication follows from noting that
$\mu' = \average(C') = \mu+x$,
$\max\braced{\precision(C'),\precision(\mu')}\le d$, and by applying the
above argument to $-x$ instead of $x$.
\end{proof}

\begin{observation}%--------
\label{obs: power of 2 rescaling}
Let $\mu = \average(C)$, $\delta = \max\braced{\precision(C),\precision(\mu)}$, 
$C' = C\cdot 2^\delta$ with $\mu'=\average(C')= 2^\delta\mu$. 
	(Thus $C'\cup\braced{\mu'}\ingroundset\integers$.)
	Then $C$ is perfectly mixable with precision $d\ge\delta$
	if and only if $C'$ is perfectly mixable with precision $d' = d-\delta$. 
\end{observation}%----------


\begin{proof}
$(\Rightarrow)$
Let $G$ be a perfect-mixing sequence for $C$, with precision $d$.
Run the same sequence $G$ on input $C'$.
If some node in $G$ produces a value $z$ on input $C$, then its value on input $C'$ will be
$z 2^\delta$, and $\precision(z 2^\delta) = \max\braced{\precision(z)-\delta,0} \le d-\delta = d'$.

$(\Leftarrow)$
Let $G'$ be a perfect-mixing sequence for $C'$ with precision $d'$.
Run $G'$ on input $C$. If some node in $G'$ produces a value $y$ on input $C'$, then
its value on input $C$ will be $y/2^\delta$, and
$\precision(y/2^\delta) \leq \precision(y)+\delta \le d'+\delta = d$.
\end{proof}


%%%%%%%

\myparagraph{Integral configurations}
Per Observation~\ref{obs: power of 2 rescaling}, we can restrict our attention to configurations with integer
values and average, that is, we will be assuming that $C\cup\braced{\mu}\ingroundset \integers$.

For $x\in\posintegers$, if each $c\in C$ is a multiple of $x$,
let $C/x = \braced{c/x \mid c\in C}$.
For integral configurations, we can extend Observation~\ref{obs: power of 2 rescaling}
to also multiplying $C$ by an odd integer or dividing it by a common odd factor of
all concentrations in $C$. 

\begin{observation}%--------
\label{obs: odd integer rescaling}
Assume that $C\cup\braced{\mu}\ingroundset \integers$ and let $x\in\posintegers$ be odd.
%
\begin{description}
%
\item{\emph{(a)}} Let $C' = C\cdot x$. Then
	$C$ is perfectly mixable with precision $0$ if and only if
	$C'$ is perfectly mixable with precision $0$.	
%
\item{\emph{(b)}} Suppose that $x$ 
is a divisor of all concentrations in $C\cup\braced{\mu}$.
Then $C$ is perfectly mixable with precision $0$ if and only if
$C/x$ is perfectly mixable with precision $0$.
\end{description}	
\end{observation}%----------

\begin{proof}
Part~(b) follows from~(a), so we only prove part~(a).
Any sequence $G$ of mixing operations for $C$ can be applied to $C\cdot x$. 
By simple induction,
if some intermediate value in $G$ was an integer $z$, now its value will be $z x$,
also an integer. This shows the $(\Rightarrow)$ implication.

To justify the $(\Leftarrow)$ implication, 
suppose that $G'$ is a sequence of mixing operations for $C'$ and that all
concentrations in $G'$ are integer. 
Since $x$ is odd and all concentrations in $C'$ are multiples of $x$, 
every concentration in $G'$ will be also a multiple of $x$, including  $\average(C')$.  
Thus, if  we run $G'$ on $C$ instead of $C'$, if some node's concentration
was $c x$, now it will be $c$.
Thus, the $(\Leftarrow)$ implication holds.
\end{proof}

%%%%%%%








Let $\binrationals$ be the set of  binary numbers.
For $c\in\binrationals$, we denote by $\precision(c)$ the precision of $c$, that is the
number of fractional bits in the binary representation of $c$, assuming there are no trailing $0$'s.
In other words, $\precision(c)$ is the smallest $d\in\nonnegintegers$ such that $c = a/2^d$ for some $a\in\integers$. 
If $c = a/2^d$ represents actual fluid concentration, then we have $0\le a \le 2^d$. However,  
it is convenient to relax this restriction and allow ``concentration values'' that are
arbitrary binary numbers, even negative. In fact, as we show shortly, 
it will be convenient to work with integral values.

By a \emph{configuration} we mean a multiset of $n$ binary numbers, called droplets or concentrations.
In the literature, multisets are often represented by their characteristic functions (that specify
the multiplicity of each element of the ground set). In this paper we will generally use set-theoretic
terminology, with its natural interpretation. For example, for a configuration $C$ and a concentration $a$, 
$a\in C$ means that the multiplicity of $a$ in $C$ is strictly positive, while $a\notin C$ means that it's zero.
Or, $C - \braced{a} = C'$ means that the multiplicity of $a$ in $C'$ is one less than in $C$, while
other multiplicities are the same.
The number of droplets in $C$ is denoted $\nofdroplets{C}=n$, 
while the number of different concentrations is denoted $\nofconcentrations{C}=m$.
We will typically denote a configuration by
$C = \braced{f_1:c_1,f_2: c_2,...,f_m:c_m} \ingroundset \binrationals$, where each $c_i$ represents
a (different) concentration value and $f_i$ denotes the
multiplicity of $c_i$ in $C$, so that $\sum_{i=1}^m f_i = n$. 
Occasionally, if it does not lead to confusion, we may say
``droplet $c_i$'' or ``concentration $c_i$'', referring to some droplet with concentration $c_i$.
If $f_i=1$, we shorten ``$f_i:c_i$'' to just ``$c_i$''.
If $f_i = 1$ we say that droplet $c_i$ is a \emph{singleton}, 
if $f_i = 2$ we say that droplet $c_i$ is a \emph{doubleton}
and if $f_i\geq 2$ we say that droplet $c_i$ is a \emph{non-singleton}.
By $\mysum(C)$ we denote the sum of $C$, that is $\mysum(C) = \sum_{c\in C} c$.
$\average(C) = \mysum(C)/n$ is the average value of the concentrations in $C$ and will
be typically denoted by $\mu$.
(Later, we will typically deal with configurations $C$ such that
$C\cup\braced{\mu}\ingroundset\integers$.)

Mixing graphs were defined in the introduction. As we are not concerned in this paper with the
topological properties of mixing graphs, we will often identify a mixing graph $G$ with a
corresponding \emph{mixing sequence}, which is a sequence (not necessarily unique) of mixing operations
that convert $C$ into its perfect mixture. In other words, a mixing sequence is a sequence
of mixing operations in a topological ordering of a mixing graph.

Of course in a perfect-mixing graph (or sequence) $G$ for $C$, all concentrations in $G$,
including those in $C\cup\braced{\mu}$,
must have finite precision (that is, belong to $\binrationals$);
in fact, the maximum concentration in $G$ is at least $\max\braced{\precision(C),\precision(\mu)}$.
In addition to the basic question about finding a perfect-mixing graph for $C$, 
we are also interested in bounding the precision required to do so.

For $x\in\binrationals$, define multisets $C+x = \braced{c+x \mid c\in C}$,
$C-x = C+(-x)$, and $C\cdot x = \braced{c\cdot x \mid c\in C}$.
The next observation says that offsetting all values in $C$
does not affect perfect mixability, as long as the offset value's precision
does not exceed that of $C$ or $\mu$.

\begin{observation}%--------
\label{obs: offsetting does not affect reachability}
Let $\mu = \average(C)$ and $x\in\binrationals$. 
Also, let $d\in\posintegers$ be such that $d\ge \max\braced{\precision(C),\precision(\mu),\precision(x)}$.
Then $C$ is perfectly mixable with precision $d$ if and only if
$C' = C+x$ is perfectly mixable with precision $d$.
\end{observation}%----------

\begin{proof}
$(\Rightarrow)$
Suppose that $G$ is a perfect-mixing sequence for $C$ with precision $d$.
Run the same sequence $G$ on input $C'$. If some mixing step in $G$ produces a value $z$ when
the input is $C$ then on $C'$ its value is $z+x$,
and $\precision(z+x)\le \max\braced{\precision(z),\precision(x)} \le d$.
Thus the maximum precision in $G$ for input $C'$ is at most $d$. 

$(\Leftarrow)$
The proof for this implication follows from noting that
$\mu' = \average(C') = \mu+x$,
$\max\braced{\precision(C'),\precision(\mu')}\le d$, and by applying the
above argument to $-x$ instead of $x$.
\end{proof}

\begin{observation}%--------
\label{obs: power of 2 rescaling}
Let $\mu = \average(C)$, $\delta = \max\braced{\precision(C),\precision(\mu)}$, 
$C' = C\cdot 2^\delta$ with $\mu'=\average(C')= 2^\delta\mu$. 
	(Thus $C'\cup\braced{\mu'}\ingroundset\integers$.)
	Then $C$ is perfectly mixable with precision $d\ge\delta$
	if and only if $C'$ is perfectly mixable with precision $d' = d-\delta$. 
\end{observation}%----------


\begin{proof}
$(\Rightarrow)$
Let $G$ be a perfect-mixing sequence for $C$, with precision $d$.
Run the same sequence $G$ on input $C'$.
If some node in $G$ produces a value $z$ on input $C$, then its value on input $C'$ will be
$z 2^\delta$, and $\precision(z 2^\delta) = \max\braced{\precision(z)-\delta,0} \le d-\delta = d'$.

$(\Leftarrow)$
Let $G'$ be a perfect-mixing sequence for $C'$ with precision $d'$.
Run $G'$ on input $C$. If some node in $G'$ produces a value $y$ on input $C'$, then
its value on input $C$ will be $y/2^\delta$, and
$\precision(y/2^\delta) \leq \precision(y)+\delta \le d'+\delta = d$.
\end{proof}


%%%%%%%

\myparagraph{Integral configurations}
Per Observation~\ref{obs: power of 2 rescaling}, we can restrict our attention to configurations with integer
values and average, that is, we will be assuming that $C\cup\braced{\mu}\ingroundset \integers$.

For $x\in\posintegers$, if each $c\in C$ is a multiple of $x$,
let $C/x = \braced{c/x \mid c\in C}$.
For integral configurations, we can extend Observation~\ref{obs: power of 2 rescaling}
to also multiplying $C$ by an odd integer or dividing it by a common odd factor of
all concentrations in $C$. 

\begin{observation}%--------
\label{obs: odd integer rescaling}
Assume that $C\cup\braced{\mu}\ingroundset \integers$ and let $x\in\posintegers$ be odd.
%
\begin{description}
%
\item{\emph{(a)}} Let $C' = C\cdot x$. Then
	$C$ is perfectly mixable with precision $0$ if and only if
	$C'$ is perfectly mixable with precision $0$.	
%
\item{\emph{(b)}} Suppose that $x$ 
is a divisor of all concentrations in $C\cup\braced{\mu}$.
Then $C$ is perfectly mixable with precision $0$ if and only if
$C/x$ is perfectly mixable with precision $0$.
\end{description}	
\end{observation}%----------

\begin{proof}
Part~(b) follows from~(a), so we only prove part~(a).
Any sequence $G$ of mixing operations for $C$ can be applied to $C\cdot x$. 
By simple induction,
if some intermediate value in $G$ was an integer $z$, now its value will be $z x$,
also an integer. This shows the $(\Rightarrow)$ implication.

To justify the $(\Leftarrow)$ implication, 
suppose that $G'$ is a sequence of mixing operations for $C'$ and that all
concentrations in $G'$ are integer. 
Since $x$ is odd and all concentrations in $C'$ are multiples of $x$, 
every concentration in $G'$ will be also a multiple of $x$, including  $\average(C')$.  
Thus, if  we run $G'$ on $C$ instead of $C'$, if some node's concentration
was $c x$, now it will be $c$.
Thus, the $(\Leftarrow)$ implication holds.
\end{proof}

%%%%%%%








Let $\binrationals$ be the set of  binary numbers.
For $c\in\binrationals$, we denote by $\precision(c)$ the precision of $c$, that is the
number of fractional bits in the binary representation of $c$, assuming there are no trailing $0$'s.
In other words, $\precision(c)$ is the smallest $d\in\nonnegintegers$ such that $c = a/2^d$ for some $a\in\integers$. 
If $c = a/2^d$ represents actual fluid concentration, then we have $0\le a \le 2^d$. However,  
it is convenient to relax this restriction and allow ``concentration values'' that are
arbitrary binary numbers, even negative. In fact, as we show shortly, 
it will be convenient to work with integral values.

By a \emph{configuration} we mean a multiset of $n$ binary numbers, called droplets or concentrations.
In the literature, multisets are often represented by their characteristic functions (that specify
the multiplicity of each element of the ground set). In this paper we will generally use set-theoretic
terminology, with its natural interpretation. For example, for a configuration $C$ and a concentration $a$, 
$a\in C$ means that the multiplicity of $a$ in $C$ is strictly positive, while $a\notin C$ means that it's zero.
Or, $C - \braced{a} = C'$ means that the multiplicity of $a$ in $C'$ is one less than in $C$, while
other multiplicities are the same.
The number of droplets in $C$ is denoted $\nofdroplets{C}=n$, 
while the number of different concentrations is denoted $\nofconcentrations{C}=m$.
We will typically denote a configuration by
$C = \braced{f_1:c_1,f_2: c_2,...,f_m:c_m} \ingroundset \binrationals$, where each $c_i$ represents
a (different) concentration value and $f_i$ denotes the
multiplicity of $c_i$ in $C$, so that $\sum_{i=1}^m f_i = n$. 
Occasionally, if it does not lead to confusion, we may say
``droplet $c_i$'' or ``concentration $c_i$'', referring to some droplet with concentration $c_i$.
If $f_i=1$, we shorten ``$f_i:c_i$'' to just ``$c_i$''.
If $f_i = 1$ we say that droplet $c_i$ is a \emph{singleton}, 
if $f_i = 2$ we say that droplet $c_i$ is a \emph{doubleton}
and if $f_i\geq 2$ we say that droplet $c_i$ is a \emph{non-singleton}.
By $\mysum(C)$ we denote the sum of $C$, that is $\mysum(C) = \sum_{c\in C} c$.
$\average(C) = \mysum(C)/n$ is the average value of the concentrations in $C$ and will
be typically denoted by $\mu$.
(Later, we will typically deal with configurations $C$ such that
$C\cup\braced{\mu}\ingroundset\integers$.)

Mixing graphs were defined in the introduction. As we are not concerned in this paper with the
topological properties of mixing graphs, we will often identify a mixing graph $G$ with a
corresponding \emph{mixing sequence}, which is a sequence (not necessarily unique) of mixing operations
that convert $C$ into its perfect mixture. In other words, a mixing sequence is a sequence
of mixing operations in a topological ordering of a mixing graph.

Of course in a perfect-mixing graph (or sequence) $G$ for $C$, all concentrations in $G$,
including those in $C\cup\braced{\mu}$,
must have finite precision (that is, belong to $\binrationals$);
in fact, the maximum concentration in $G$ is at least $\max\braced{\precision(C),\precision(\mu)}$.
In addition to the basic question about finding a perfect-mixing graph for $C$, 
we are also interested in bounding the precision required to do so.

For $x\in\binrationals$, define multisets $C+x = \braced{c+x \mid c\in C}$,
$C-x = C+(-x)$, and $C\cdot x = \braced{c\cdot x \mid c\in C}$.
The next observation says that offsetting all values in $C$
does not affect perfect mixability, as long as the offset value's precision
does not exceed that of $C$ or $\mu$.

\begin{observation}%--------
\label{obs: offsetting does not affect reachability}
Let $\mu = \average(C)$ and $x\in\binrationals$. 
Also, let $d\in\posintegers$ be such that $d\ge \max\braced{\precision(C),\precision(\mu),\precision(x)}$.
Then $C$ is perfectly mixable with precision $d$ if and only if
$C' = C+x$ is perfectly mixable with precision $d$.
\end{observation}%----------

\begin{proof}
$(\Rightarrow)$
Suppose that $G$ is a perfect-mixing sequence for $C$ with precision $d$.
Run the same sequence $G$ on input $C'$. If some mixing step in $G$ produces a value $z$ when
the input is $C$ then on $C'$ its value is $z+x$,
and $\precision(z+x)\le \max\braced{\precision(z),\precision(x)} \le d$.
Thus the maximum precision in $G$ for input $C'$ is at most $d$. 

$(\Leftarrow)$
The proof for this implication follows from noting that
$\mu' = \average(C') = \mu+x$,
$\max\braced{\precision(C'),\precision(\mu')}\le d$, and by applying the
above argument to $-x$ instead of $x$.
\end{proof}

\begin{observation}%--------
\label{obs: power of 2 rescaling}
Let $\mu = \average(C)$, $\delta = \max\braced{\precision(C),\precision(\mu)}$, 
$C' = C\cdot 2^\delta$ with $\mu'=\average(C')= 2^\delta\mu$. 
	(Thus $C'\cup\braced{\mu'}\ingroundset\integers$.)
	Then $C$ is perfectly mixable with precision $d\ge\delta$
	if and only if $C'$ is perfectly mixable with precision $d' = d-\delta$. 
\end{observation}%----------


\begin{proof}
$(\Rightarrow)$
Let $G$ be a perfect-mixing sequence for $C$, with precision $d$.
Run the same sequence $G$ on input $C'$.
If some node in $G$ produces a value $z$ on input $C$, then its value on input $C'$ will be
$z 2^\delta$, and $\precision(z 2^\delta) = \max\braced{\precision(z)-\delta,0} \le d-\delta = d'$.

$(\Leftarrow)$
Let $G'$ be a perfect-mixing sequence for $C'$ with precision $d'$.
Run $G'$ on input $C$. If some node in $G'$ produces a value $y$ on input $C'$, then
its value on input $C$ will be $y/2^\delta$, and
$\precision(y/2^\delta) \leq \precision(y)+\delta \le d'+\delta = d$.
\end{proof}


%%%%%%%

\myparagraph{Integral configurations}
Per Observation~\ref{obs: power of 2 rescaling}, we can restrict our attention to configurations with integer
values and average, that is, we will be assuming that $C\cup\braced{\mu}\ingroundset \integers$.

For $x\in\posintegers$, if each $c\in C$ is a multiple of $x$,
let $C/x = \braced{c/x \mid c\in C}$.
For integral configurations, we can extend Observation~\ref{obs: power of 2 rescaling}
to also multiplying $C$ by an odd integer or dividing it by a common odd factor of
all concentrations in $C$. 

\begin{observation}%--------
\label{obs: odd integer rescaling}
Assume that $C\cup\braced{\mu}\ingroundset \integers$ and let $x\in\posintegers$ be odd.
%
\begin{description}
%
\item{\emph{(a)}} Let $C' = C\cdot x$. Then
	$C$ is perfectly mixable with precision $0$ if and only if
	$C'$ is perfectly mixable with precision $0$.	
%
\item{\emph{(b)}} Suppose that $x$ 
is a divisor of all concentrations in $C\cup\braced{\mu}$.
Then $C$ is perfectly mixable with precision $0$ if and only if
$C/x$ is perfectly mixable with precision $0$.
\end{description}	
\end{observation}%----------

\begin{proof}
Part~(b) follows from~(a), so we only prove part~(a).
Any sequence $G$ of mixing operations for $C$ can be applied to $C\cdot x$. 
By simple induction,
if some intermediate value in $G$ was an integer $z$, now its value will be $z x$,
also an integer. This shows the $(\Rightarrow)$ implication.

To justify the $(\Leftarrow)$ implication, 
suppose that $G'$ is a sequence of mixing operations for $C'$ and that all
concentrations in $G'$ are integer. 
Since $x$ is odd and all concentrations in $C'$ are multiples of $x$, 
every concentration in $G'$ will be also a multiple of $x$, including  $\average(C')$.  
Thus, if  we run $G'$ on $C$ instead of $C'$, if some node's concentration
was $c x$, now it will be $c$.
Thus, the $(\Leftarrow)$ implication holds.
\end{proof}

%%%%%%%






